\section{Problem Setup}
\label{sec:setup}

This section establishes the experimental framework for our mechanism-driven empirical study. We clarify the research objective, the branch-wise KD structure, the visibility degradation protocol, and the evaluation methodology.

\subsection{Research Objective}

Unlike conventional KD research that proposes new architectures or loss functions, our goal is \textbf{mechanism identification}: determining \emph{why} KD fails under visibility degradation. We seek to answer:
\begin{enumerate}
    \item Does KD failure exhibit branch-wise structure, or do all branches degrade uniformly?
    \item Which hypothesized mechanisms (distribution mismatch, uncertainty amplification, visibility disruption) receive empirical support?
    \item What implications do these findings have for future method design?
\end{enumerate}

\subsection{Branch-Wise KD Framework}

We adopt a modular KD framework that isolates four distinct knowledge transfer mechanisms:

\begin{itemize}
    \item \textbf{Logit distillation ($\mathcal{L}_{\text{logit}}$)}: Matches softened classification outputs between teacher and student using temperature-scaled cross-entropy \cite{hinton2015distilling}:
    \begin{equation}
    \mathcal{L}_{\text{logit}} = -\sum_{i} p_i^T \log(q_i^S)
    \end{equation}
    where $p_i^T$ and $q_i^S$ are teacher and student probability distributions with temperature $T$.

    \item \textbf{Feature distillation ($\mathcal{L}_{\text{feat}}$)}: Aligns intermediate feature representations through $L_2$ distance or cosine similarity \cite{romero2014fitnets}.

    \item \textbf{Attention distillation ($\mathcal{L}_{\text{attn}}$)}: Transfers spatial attention maps derived from feature activations \cite{zagoruyko2016paying}.

    \item \textbf{Localization distillation ($\mathcal{L}_{\text{loc}}$)}: Specifically distills bounding box regression outputs, including coordinate and dimension predictions \cite{dai2021general}.
\end{itemize}

We also include a \textbf{student-only} baseline trained without any teacher supervision, enabling measurement of relative KD gain for each branch.

\subsection{Visibility Degradation Protocol}

To simulate realistic visibility degradation, we employ the standard atmospheric scattering model used in Foggy Cityscapes \cite{sakaridis2018semantic}:

\begin{equation}
I(x) = J(x) \cdot t(x) + A \cdot (1 - t(x))
\end{equation}

where $I(x)$ is the observed foggy image, $J(x)$ is the scene radiance, $A$ is the global atmospheric light, and $t(x) = e^{-\beta \cdot d(x)}$ is the transmission map with scattering coefficient $\beta$ and depth $d(x)$.

We define three visibility levels by varying $\beta$:
\begin{itemize}
    \item \textbf{Light}: $\beta = 0.005$ (minimal visibility reduction)
    \item \textbf{Moderate}: $\beta = 0.01$ (noticeable fog, objects remain identifiable)
    \item \textbf{Heavy}: $\beta = 0.02$ (significant occlusion and contrast reduction)
\end{itemize}

\subsection{Model and Dataset Configuration}

\textbf{Teacher model}: YOLOv8l (large variant, ~25.9M parameters), trained on clear-weather Cityscapes images and frozen during student training.

\textbf{Student model}: YOLOv8m (medium variant), trained from scratch under each visibility condition.

\textbf{Dataset}: Foggy Cityscapes (YOLO format), containing 8 object categories (person, rider, car, truck, bus, train, motorcycle, bicycle) with 2,975 training and 500 validation images.

\textbf{Training configuration}: 150 epochs, batch size 32, image size $640 \times 640$, AdamW optimizer with initial learning rate 0.005 and cosine annealing. Standard augmentations (Mosaic, Mixup) applied consistently across all experiments.

\subsection{Evaluation Metrics}

\textbf{Primary metric}: mean Average Precision at IoU=0.50 (mAP@50), the standard detection metric measuring both localization accuracy and classification correctness.

\textbf{Secondary metrics}: For mechanism analysis, we compute:
\begin{itemize}
    \item \textbf{KD Gain}: $\Delta_{\text{branch}} = \text{mAP}_{\text{branch}} - \text{mAP}_{\text{student-only}}$ at each visibility level
    \item \textbf{Ranking stability}: Kendall's $\tau$ correlation between branch rankings across visibility conditions
    \item \textbf{Mechanism correlations}: Pearson $r$ between hypothesized factors (divergence, entropy, occlusion) and KD gains
\end{itemize}

\subsection{Experimental Matrix}

Our complete experimental design forms a $5 \times 3$ matrix:
\begin{itemize}
    \item 5 KD configurations: student-only, logit-only, feature-only, attention-only, localization-only
    \item 3 visibility levels: light, moderate, heavy
    \item Total: 15 training runs, all controlled for random seed and hyperparameters
\end{itemize}

This structure enables both observation-level analysis (comparing branch performance) and mechanism-level analysis (testing specific hypotheses against measured outcomes).
